

Los resultados experimentales demuestran que la superioridad asintótica teórica no siempre dicta el rendimiento práctico en instancias de tamaño moderado. En el problema de ordenamiento, la implementación de biblioteca \texttt{std::sort} probó ser la más robusta y veloz en todas las distribuciones de datos evaluadas, minimizando además el consumo de memoria; contrastando fuertemente con algoritmos como Patience Sort, cuyo tiempo se disparó en arreglos aleatorios y ascendentes. Por su parte, la multiplicación de matrices evidenció que la sobrecarga computacional de las llamadas recursivas y su mayor huella de memoria provocaron que la aproximación iterativa Naive fuera invariablemente más rápida para todos los tipos de matrices dentro del dominio estudiado (hasta $n=256$). Esto confirma empíricamente que Strassen requiere de situaciones especificas para justificar su utilidad en la práctica.